% ============================================================
%  Chapter 10 — The Future of Robots
%  Robot World: Meet the Machines That Help Us
% ============================================================

\chapter{The Future of Robots \textemdash{} Safe, Kind, and Smart}

\chapteropening{%
  Robots are getting smarter, smaller, softer, and more helpful every year.
  But with great power comes great responsibility.
  What will Robot World look like when you grow up \textemdash{} and how
  can we make sure robots stay safe, kind, and fair?
}
\chapterbadge{Tomorrow's Robots \textemdash{} Safe, Kind, and Yours to Shape!}

% --- Introduction ---
\section*{What Comes Next?}

\character{Dot}{We have met factory arms, drones, rovers, surgical bots,
  and companion robots. But the most exciting robots have not been
  invented yet \textemdash{} and one of them might be yours!}

\sectionrule

% --- Trends ---
\section*{Five Big Trends}

\subsection*{1. Soft Robots}

Instead of rigid metal, some new robots are made of \term{silicone}, rubber,
or even living cells. Inspired by octopuses and caterpillars, they can
squeeze through tiny gaps and handle delicate objects without crushing them.

\subsection*{2. Swarm Robots}

Hundreds of tiny, simple robots working together like ants. No single robot
is smart, but the \term{swarm} as a whole can build structures, search large areas,
or move heavy objects cooperatively.

\begin{funfact}
  Researchers at Harvard created ``Kilobots'' \textemdash{} swarms of over
  1,000 coin-sized robots that arrange themselves into shapes
  without any central leader.
\end{funfact}

\subsection*{3. Learning Robots}

Using \term{machine learning}, robots can improve by practising \textemdash{}
much like you get better at a video game the more you play.
A robot arm might try thousands of ways to pick up an egg until
it finds the perfect gentle grip.

\subsection*{4. Bio-Inspired Robots}

Engineers copy nature's best ideas:
\begin{itemize}
  \item \term{Cheetah}-inspired legs for fast running.
  \item \term{Gecko}-inspired feet for climbing walls.
  \item Fish-inspired tails for silent swimming.
\end{itemize}

\subsection*{5. Robots and AI}

As \term{artificial intelligence} improves, robots will understand language,
recognise more objects, and make better decisions in unfamiliar situations.

\begin{bigidea}
  The future of robots is not about replacing humans \textemdash{}
  it is about humans and robots working together, each doing
  what they do best.
\end{bigidea}

\sectionrule

% --- Jobs of the future ---
\section*{Robot Jobs of the Future}

\begin{itemize}
  \item Farming robots that pick fruit without bruising it.
  \item Construction robots that 3D-print entire houses.
  \item Ocean-cleaning robots that collect plastic from the sea.
  \item Personal tutor robots that adapt lessons to each student.
  \item Disaster-response drones that arrive before ambulances.
\end{itemize}

\begin{picturethis}
  Imagine waking up in a house that was 3D-printed by a robot overnight,
  eating breakfast picked by a farming robot that morning, and riding
  to school in a \term{self-driving} bus. Much of this technology already exists
  \textemdash{} it just needs time to become affordable and reliable.
\end{picturethis}

\sectionrule


% --- Safety and responsibility ---
\section*{Building Robots Responsibly}

Great robot designers ask more than ``Can we build it?'' They also ask:

\begin{itemize}
  \item \textbf{Is it safe?} --- Robots near people need \term{stop buttons}, soft edges,
    speed limits, and careful sensors.
  \item \textbf{Is it fair?} --- AI systems can learn unfair patterns if trained
    on biased examples. Engineers must test carefully.
  \item \textbf{Does it respect \term{privacy}?} --- A robot with a camera or microphone
    should only collect information it truly needs, and keep it secure.
  \item \textbf{Can humans stay in control?} --- Every robot should have a way for
    a person to stop it instantly.
\end{itemize}

\begin{bigidea}
  The best robots are not just clever --- they are safe, kind, useful,
  and designed with people's wellbeing in mind.
  Technology is a tool; \emph{you} decide how it is used.
\end{bigidea}

\begin{picturethis}
  Imagine designing a playground. You want it to be fun, but you also need
  strong materials, soft ground beneath the swings, fences near roads,
  and rules so everyone can enjoy it safely.
  Designing robots is the same --- fun \emph{and} responsibility together.
\end{picturethis}

\sectionrule

% --- Your role ---
\section*{Your Place in Robot World}

You do not have to be an engineer to shape the future of robots.
The world needs people who:
\begin{itemize}
  \item \textbf{Design} robots that are safe and fair.
  \item \textbf{Write rules} about where robots should and should not be used.
  \item \textbf{Create art} for robot interfaces and stories about our shared future.
  \item \textbf{Ask questions} \textemdash{} the hardest but most important job of all.
\end{itemize}

\character{Cora}{Whatever you choose to do when you grow up, robots will
  be part of your world. Understanding them now means you can help
  decide how they are used \textemdash{} for good.}

\sectionrule

% --- Try It ---
\begin{tryit}
  \textbf{Invent a robot.}
  \begin{enumerate}
    \item Choose a problem you care about (litter, loneliness, homework\ldots).
    \item Sketch a robot that could help solve it.
    \item Label its sensors, actuators, and brain.
    \item Write three steps of its main algorithm.
    \item Give it a name!
  \end{enumerate}
  Who knows \textemdash{} maybe one day you will build it for real.
\end{tryit}

\sectionrule

% --- Closing ---
\section*{Goodbye from Dot}

\character{Dot}{Thank you for exploring Robot World with me!
  We have met factory arms, flying drones, deep-sea divers, surgical
  micro-bots, and companion friends. Remember: every robot starts as
  an idea in someone's head. The next great robot might start in yours.
  See you in the next adventure!}

\sectionrule


\begin{quickcheck}
  \begin{enumerate}
    \item What are swarm robots and why are they useful?
    \item Name one way engineers copy nature when designing robots.
    \item Why should robot designers think about safety and privacy?
  \end{enumerate}
\end{quickcheck}

\sectionrule
% --- New Words ---
\begin{newwords}
  \begin{description}[labelwidth=3.8cm, leftmargin=4.0cm]
    \item[Soft robot] A robot made of flexible materials that can bend and stretch.
    \item[Swarm robotics] Many simple robots cooperating to accomplish complex tasks.
    \item[Machine learning] A type of AI where computers improve through experience.
    \item[Bio-inspired] Designed by copying ideas from nature.
    \item[3D printing] Building objects layer by layer from a digital model.
    \item[Artificial intelligence] Computers that can perform tasks normally requiring human thought.
  \end{description}
\end{newwords}
